\documentclass[11pt]{article}

\usepackage{amsmath, amssymb, amsthm}
\usepackage{mathtools}
\usepackage{geometry}
\usepackage{enumitem}
\usepackage{bm}
\usepackage{tcolorbox}

\geometry{margin=1in}

\newtheorem{theorem}{Theorem}[section]
\newtheorem{proposition}[theorem]{Proposition}
\newtheorem{lemma}[theorem]{Lemma}
\newtheorem{corollary}[theorem]{Corollary}
\newtheorem{assumption}[theorem]{Assumption}
\theoremstyle{definition}
\newtheorem{definition}[theorem]{Definition}
\newtheorem{example}[theorem]{Example}
\newtheorem{remark}[theorem]{Remark}

\DeclareMathOperator{\Tr}{Tr}
\DeclareMathOperator{\rank}{rank}
\DeclareMathOperator{\supp}{supp}
\DeclareMathOperator{\Range}{Range}
\DeclareMathOperator{\argmax}{arg\,max}
\DeclareMathOperator{\argmin}{arg\,min}
\DeclareMathOperator{\TV}{TV}
\DeclareMathOperator{\KL}{KL}
\DeclareMathOperator{\Hell}{H}
\DeclareMathOperator{\diam}{diam}

\newcommand{\R}{\mathbb{R}}
\newcommand{\N}{\mathbb{N}}
\newcommand{\E}{\mathbb{E}}
\newcommand{\Prob}{\mathbb{P}}
\newcommand{\Splus}{\mathcal{S}_+^d}
\newcommand{\Splusplus}{\mathcal{S}_{++}^d}
\newcommand{\Srank}{\mathcal{S}_{\leq r}^d}
\newcommand{\HK}{\mathcal{H}_K}
\newcommand{\X}{\mathcal{X}}
\newcommand{\M}{\mathcal{M}}
\newcommand{\GP}{\mathcal{GP}}
\newcommand{\iid}{\stackrel{\mathrm{iid}}{\sim}}
\newcommand{\inprod}[2]{\langle #1, #2 \rangle}
\newcommand{\norm}[1]{\| #1 \|}
\newcommand{\op}{\mathrm{op}}

\title{A General Rate Theorem for Mixing Measure Estimation}
\author{}
\date{}

\begin{document}

\maketitle

%%%%%%%%%%%%%%%%%%%%%%%%%%%%%%%%%%%%%%%%%%%%%%%%%%%%%%%%%%%%%%%%%%%%%%%%%%%%%%%
\section{The Meta-Theorem}
%%%%%%%%%%%%%%%%%%%%%%%%%%%%%%%%%%%%%%%%%%%%%%%%%%%%%%%%%%%%%%%%%%%%%%%%%%%%%%%

\begin{tcolorbox}[colback=blue!5!white, colframe=blue!50!black, title=Main Result]
\begin{theorem}[General Rate Theorem]
\label{thm:main}
Let $\theta^* \in \Theta \subseteq \M_+(\Splus)$. Suppose the model $\{K_\theta : \theta \in \Theta\}$ has \textbf{local metric entropy} around $\theta^*$ satisfying:
\begin{equation}
    \log N\left(\epsilon, \{\theta \in \Theta : d_K(\theta, \theta^*) \leq \delta\}, d_K\right) \leq D(\theta^*) \cdot \log\left(\frac{\delta}{\epsilon}\right)
\end{equation}
for all $0 < \epsilon < \delta$ sufficiently small. Then:
\begin{equation}
    \boxed{
    d_K(\hat{\theta}_m, \theta^*) = O_p\left(m^{-\frac{1}{2 + D(\theta^*)/s}}\right)
    }
\end{equation}
where $s$ is the smoothness parameter if $\theta^*$ has a density, or $s = \infty$ (giving rate $m^{-1/2}$) in the parametric case.

More precisely, with $D = D(\theta^*)$:
\begin{equation}
    \boxed{
    d_K(\hat{\theta}_m, \theta^*) = O_p\left(m^{-\alpha(\theta^*)}\right), \quad \alpha(\theta^*) = \frac{1}{2 + D(\theta^*)}
    }
\end{equation}
when the density has regularity $s = 1$.
\end{theorem}
\end{tcolorbox}

\begin{definition}[Effective Dimension]
\label{def:eff-dim}
The \textbf{effective dimension} of $\theta^*$ relative to the model $\Theta$ is:
\begin{equation}
    D(\theta^*) := \lim_{\epsilon \to 0} \frac{\log N(\epsilon, \Theta_{\theta^*}, d_K)}{\log(1/\epsilon)}
\end{equation}
where $\Theta_{\theta^*}$ is a local neighborhood of $\theta^*$ in $\Theta$, and $N(\epsilon, \cdot, d_K)$ is the $\epsilon$-covering number in the kernel metric $d_K(\theta_1, \theta_2) = \|K_{\theta_1} - K_{\theta_2}\|$.
\end{definition}

\begin{remark}
The effective dimension $D(\theta^*)$ captures the \textbf{local complexity} of the model around $\theta^*$. It can vary with $\theta^*$---this is the key to adaptation.
\end{remark}

%%%%%%%%%%%%%%%%%%%%%%%%%%%%%%%%%%%%%%%%%%%%%%%%%%%%%%%%%%%%%%%%%%%%%%%%%%%%%%%
\section{The Exponent Formula}
%%%%%%%%%%%%%%%%%%%%%%%%%%%%%%%%%%%%%%%%%%%%%%%%%%%%%%%%%%%%%%%%%%%%%%%%%%%%%%%

\begin{tcolorbox}[colback=green!5!white, colframe=green!50!black, title=Rate Exponent]
For $\theta^*$ with regularity $s$ (number of derivatives of density, or $s = \infty$ for parametric):
\begin{equation}
    \boxed{
    \alpha(\theta^*) = \frac{s}{2s + D(\theta^*)}
    }
\end{equation}
giving rate $d_K(\hat{\theta}_m, \theta^*) = O_p(m^{-\alpha(\theta^*)})$.
\end{tcolorbox}

\begin{theorem}[Explicit Exponent Formula]
\label{thm:exponent}
Let $\theta^*$ belong to a class $\Theta$ characterized by:
\begin{itemize}
    \item \textbf{Intrinsic dimension:} $\dim(\supp(\theta^*)) = D_{\mathrm{int}}$
    \item \textbf{Smoothness:} density $p^* = d\theta^*/d\mu$ has $s$ derivatives
    \item \textbf{Sparsity:} $\theta^*$ is supported on $J$ atoms (if discrete)
\end{itemize}
Then:
\begin{equation}
    \alpha(\theta^*) = \begin{cases}
        \displaystyle \frac{1}{2} & \text{if } \theta^* \text{ is } J\text{-sparse (parametric)} \\[1.5em]
        \displaystyle \frac{s}{2s + D_{\mathrm{int}}} & \text{if } \theta^* \text{ has } C^s \text{ density on } D_{\mathrm{int}}\text{-dim support}
    \end{cases}
\end{equation}
\end{theorem}

%%%%%%%%%%%%%%%%%%%%%%%%%%%%%%%%%%%%%%%%%%%%%%%%%%%%%%%%%%%%%%%%%%%%%%%%%%%%%%%
\section{Recovering All Special Cases}
%%%%%%%%%%%%%%%%%%%%%%%%%%%%%%%%%%%%%%%%%%%%%%%%%%%%%%%%%%%%%%%%%%%%%%%%%%%%%%%

\begin{proposition}[Unification]
All previously stated rates are special cases of Theorem~\ref{thm:main}:

\begin{center}
\renewcommand{\arraystretch}{1.6}
\begin{tabular}{|l|c|c|c|}
\hline
\textbf{Structure} & $D(\theta^*)$ & $\alpha = \frac{s}{2s + D}$ & \textbf{Rate} \\
\hline\hline
$J$-sparse & $0$ (parametric) & $\frac{1}{2}$ & $m^{-1/2} \cdot \sqrt{Jd^2}$ \\
\hline
$C^s$ on $\Splus$ & $\frac{d(d+1)}{2s}$ & $\frac{s}{2s + d(d+1)/2}$ & $m^{-\frac{s}{2s + d(d+1)/2}}$ \\
\hline
$C^s$ on rank-$r$ & $\frac{rd}{s}$ & $\frac{s}{2s + rd}$ & $m^{-\frac{s}{2s + rd}}$ \\
\hline
$C^s$ on rank-$1$ & $\frac{d}{s}$ & $\frac{s}{2s + d}$ & $m^{-\frac{s}{2s + d}}$ \\
\hline
$(\epsilon, J)$-approx.\ sparse & $\frac{1}{\alpha}$ & $\alpha$ & $m^{-\alpha}$ \\
\hline
\end{tabular}
\end{center}
\end{proposition}

\begin{proof}
We verify each case:

\textbf{Case 1: $J$-sparse.} The parameter space is $\R^J_+ \times (\Splusplus)^J$ with dimension $p = J(1 + d(d+1)/2) \approx Jd^2$. Local covering number: $N(\epsilon, B_\delta, d_K) \lesssim (\delta/\epsilon)^p$, so $D = 0$ in the sense that $\log N \sim p \log(1/\epsilon)$ is constant in the exponent. The rate is $m^{-1/2} \cdot \sqrt{p} = m^{-1/2} \cdot \sqrt{Jd^2}$.

\textbf{Case 2: $C^s$ density on $\Splus$.} The intrinsic dimension is $D_{\mathrm{int}} = d(d+1)/2$. By standard metric entropy for $C^s$ functions on $D_{\mathrm{int}}$-dimensional domains: $\log N(\epsilon) \sim \epsilon^{-D_{\mathrm{int}}/s}$. Thus $D(\theta^*) = D_{\mathrm{int}}/s$ and $\alpha = s/(2s + D_{\mathrm{int}})$.

\textbf{Case 3: Rank-$r$ support.} The manifold $\mathcal{S}_{\leq r}^d$ has dimension $rd - \binom{r}{2} + r \approx rd$. So $D_{\mathrm{int}} = rd$ and $\alpha = s/(2s + rd)$.

\textbf{Case 4: Approximate sparsity with $\alpha$-decay.} If $\inf_{\nu \in \Theta_J} d_K(\theta^*, \nu) \lesssim J^{-\alpha}$, then balancing bias $J^{-\alpha}$ against variance $\sqrt{Jd^2/m}$ gives optimal $J \asymp m^{1/(2\alpha+1)}$ and rate $m^{-\alpha/(2\alpha+1)}$. This corresponds to effective dimension $D = 1/\alpha - 2$.
\end{proof}

%%%%%%%%%%%%%%%%%%%%%%%%%%%%%%%%%%%%%%%%%%%%%%%%%%%%%%%%%%%%%%%%%%%%%%%%%%%%%%%
\section{The General Statement}
%%%%%%%%%%%%%%%%%%%%%%%%%%%%%%%%%%%%%%%%%%%%%%%%%%%%%%%%%%%%%%%%%%%%%%%%%%%%%%%

\begin{tcolorbox}[colback=red!5!white, colframe=red!50!black, title=Most General Form]
\begin{theorem}[Master Theorem]
\label{thm:master}
Let $\theta^* \in \M_+(\Splus)$ and define:
\begin{equation}
    D(\theta^*) := \inf\left\{D \geq 0 : \exists \, \Theta \ni \theta^* \text{ with } \log N(\epsilon, \Theta, d_K) \lesssim \epsilon^{-D}\right\}
\end{equation}
the \textbf{metric entropy dimension} of the smallest model containing $\theta^*$.

Then the minimax rate over any class $\Theta \ni \theta^*$ with entropy dimension $D$ is:
\begin{equation}
    \boxed{
    \inf_{\hat{\theta}} \sup_{\theta \in \Theta} \E\left[d_K(\hat{\theta}, \theta)^2\right] \asymp m^{-\frac{2}{2 + D}}
    }
\end{equation}

If $\theta^*$ has a density with $s$ derivatives, replace $D$ with $D/s$:
\begin{equation}
    \boxed{
    \text{Rate} = m^{-\frac{2s}{2s + D_{\mathrm{int}}}}
    }
\end{equation}
where $D_{\mathrm{int}} = \dim(\supp(\theta^*))$.
\end{theorem}
\end{tcolorbox}

\begin{remark}[Interpretation]
The exponent $\alpha(\theta^*) = s/(2s + D_{\mathrm{int}})$ interpolates between:
\begin{itemize}
    \item $\alpha = 1/2$ when $D_{\mathrm{int}} = 0$ (parametric/sparse)
    \item $\alpha \to 0$ as $D_{\mathrm{int}} \to \infty$ (curse of dimensionality)
    \item $\alpha \to 1/2$ as $s \to \infty$ (infinite smoothness recovers parametric rate)
\end{itemize}
\end{remark}

%%%%%%%%%%%%%%%%%%%%%%%%%%%%%%%%%%%%%%%%%%%%%%%%%%%%%%%%%%%%%%%%%%%%%%%%%%%%%%%
\section{Adaptive Estimation}
%%%%%%%%%%%%%%%%%%%%%%%%%%%%%%%%%%%%%%%%%%%%%%%%%%%%%%%%%%%%%%%%%%%%%%%%%%%%%%%

\begin{theorem}[Adaptive Rate]
\label{thm:adaptive}
Consider the penalized estimator:
\begin{equation}
    \hat{\theta}_m = \argmin_{\theta \in \bigcup_{k=1}^\infty \Theta_k} \left\{-\ell_m(\theta) + \mathrm{pen}(\Theta_k)\right\}
\end{equation}
where $\{\Theta_k\}$ is a nested sequence of models with increasing complexity (e.g., $\Theta_k = $ measures supported on $\leq k$ atoms).

If $\mathrm{pen}(\Theta_k) \asymp D_k \log m$ where $D_k = $ effective dimension of $\Theta_k$, then:
\begin{equation}
    d_K(\hat{\theta}_m, \theta^*) = O_p\left(m^{-\alpha(\theta^*)} (\log m)^{\beta}\right)
\end{equation}
where $\alpha(\theta^*)$ is the \textbf{oracle rate} for the smallest model containing $\theta^*$, achieved \textbf{without knowing} $D(\theta^*)$ or $s$.
\end{theorem}

\begin{remark}
This is \textbf{adaptation}: the estimator automatically achieves the rate corresponding to the true complexity of $\theta^*$, paying only a logarithmic price.
\end{remark}

%%%%%%%%%%%%%%%%%%%%%%%%%%%%%%%%%%%%%%%%%%%%%%%%%%%%%%%%%%%%%%%%%%%%%%%%%%%%%%%
\section{Equivalent Characterizations of $D(\theta^*)$}
%%%%%%%%%%%%%%%%%%%%%%%%%%%%%%%%%%%%%%%%%%%%%%%%%%%%%%%%%%%%%%%%%%%%%%%%%%%%%%%

The effective dimension $D(\theta^*)$ admits several equivalent characterizations:

\begin{theorem}[Equivalent Definitions]
\label{thm:equiv}
The following are equivalent (up to constants):
\begin{enumerate}[label=(\roman*)]
    \item \textbf{Metric entropy:}
    \begin{equation}
        D = \lim_{\epsilon \to 0} \frac{\log N(\epsilon, \Theta, d_K)}{\log(1/\epsilon)}
    \end{equation}
    
    \item \textbf{Gaussian complexity:}
    \begin{equation}
        D = \lim_{\epsilon \to 0} \frac{\log \E\left[\sup_{\theta : d_K(\theta, \theta^*) \leq \epsilon} G_\theta\right]^2}{\log(1/\epsilon)}
    \end{equation}
    where $G_\theta$ is a Gaussian process indexed by $\Theta$.
    
    \item \textbf{Local Rademacher complexity:}
    \begin{equation}
        D = \lim_{\epsilon \to 0} \frac{\log \mathcal{R}_m(\{\theta : d_K(\theta, \theta^*) \leq \epsilon\})}{\log(1/\epsilon)}
    \end{equation}
    
    \item \textbf{Minkowski dimension:}
    \begin{equation}
        D = \dim_{\mathrm{Mink}}(\supp(\theta^*)) / s
    \end{equation}
    when $\theta^*$ has a $C^s$ density on its support.
    
    \item \textbf{Approximation-theoretic:}
    \begin{equation}
        D = \inf\left\{D' : \inf_{\nu \in \Theta_J} d_K(\theta^*, \nu) \lesssim J^{-1/D'} \text{ for finite mixtures } \Theta_J\right\}
    \end{equation}
\end{enumerate}
\end{theorem}

%%%%%%%%%%%%%%%%%%%%%%%%%%%%%%%%%%%%%%%%%%%%%%%%%%%%%%%%%%%%%%%%%%%%%%%%%%%%%%%
\section{Information-Theoretic Formulation}
%%%%%%%%%%%%%%%%%%%%%%%%%%%%%%%%%%%%%%%%%%%%%%%%%%%%%%%%%%%%%%%%%%%%%%%%%%%%%%%

\begin{theorem}[Information-Theoretic Rate]
\label{thm:info}
Let $I(\theta^*; \epsilon)$ denote the mutual information between $\theta$ and $(f_1, \ldots, f_m)$ when $\theta$ is uniform on the $\epsilon$-ball around $\theta^*$. The minimax rate $\epsilon_m^*$ satisfies:
\begin{equation}
    m \cdot I(\theta^*; \epsilon_m^*) \asymp \log N(\epsilon_m^*, \Theta, d_K)
\end{equation}
Solving for $\epsilon_m^*$ gives:
\begin{equation}
    \epsilon_m^* \asymp m^{-\frac{1}{2 + D(\theta^*)}}
\end{equation}
\end{theorem}

\begin{remark}
This is the \textbf{Yang-Barron} (1999) information-theoretic characterization: the rate is determined by balancing the information gained per sample against the entropy of the parameter space.
\end{remark}

%%%%%%%%%%%%%%%%%%%%%%%%%%%%%%%%%%%%%%%%%%%%%%%%%%%%%%%%%%%%%%%%%%%%%%%%%%%%%%%
\section{Summary: The Rate Formula}
%%%%%%%%%%%%%%%%%%%%%%%%%%%%%%%%%%%%%%%%%%%%%%%%%%%%%%%%%%%%%%%%%%%%%%%%%%%%%%%

\begin{tcolorbox}[colback=yellow!10!white, colframe=orange!50!black, title=\textbf{The Universal Rate Formula}]

For $\theta^* \in \M_+(\Splus)$, define:
\begin{align}
    D_{\mathrm{int}}(\theta^*) &= \text{intrinsic dimension of } \supp(\theta^*) \\
    s(\theta^*) &= \text{smoothness (derivatives) of density } d\theta^*/d\mu \\
    J(\theta^*) &= \text{number of atoms if discrete}
\end{align}

\textbf{The rate exponent is:}
\begin{equation}
    \boxed{
    \alpha(\theta^*) = \begin{cases}
        \dfrac{1}{2} & \theta^* \text{ is } J\text{-sparse} \\[1em]
        \dfrac{s}{2s + D_{\mathrm{int}}} & \theta^* \text{ has } C^s \text{ density}
    \end{cases}
    }
\end{equation}

\textbf{The rate is:}
\begin{equation}
    \boxed{
    d_K(\hat{\theta}_m, \theta^*) = O_p\left(m^{-\alpha(\theta^*)}\right) \quad \text{(up to } \sqrt{\cdot} \text{ and } \log \text{ factors)}
    }
\end{equation}

\textbf{The prefactor depends on:}
\begin{itemize}
    \item $J d^2$ for sparse (number of parameters)
    \item $\|p^*\|_{C^s}$ for smooth (regularity constant)
    \item Separation $\Delta$ for mixtures (identifiability)
\end{itemize}
\end{tcolorbox}

%%%%%%%%%%%%%%%%%%%%%%%%%%%%%%%%%%%%%%%%%%%%%%%%%%%%%%%%%%%%%%%%%%%%%%%%%%%%%%%
\section{References}
%%%%%%%%%%%%%%%%%%%%%%%%%%%%%%%%%%%%%%%%%%%%%%%%%%%%%%%%%%%%%%%%%%%%%%%%%%%%%%%

\begin{enumerate}[label={[\arabic*]}]
    \item Y.~Yang and A.~Barron. Information-theoretic determination of minimax rates of convergence. \textit{Annals of Statistics}, 27(5):1564--1599, 1999.
    
    \item S.~Ghosal, J.K.~Ghosh, and A.W.~van der Vaart. Convergence rates of posterior distributions. \textit{Annals of Statistics}, 28(2):500--531, 2000.
    
    \item L.~Birgé. Model selection via testing: an alternative to (penalized) maximum likelihood estimators. \textit{Annales de l'IHP Probabilités et Statistiques}, 42(3):273--325, 2006.
    
    \item A.B.~Tsybakov. \textit{Introduction to Nonparametric Estimation}. Springer, 2009.
    
    \item S.~Ghosal and A.W.~van der Vaart. \textit{Fundamentals of Nonparametric Bayesian Inference}. Cambridge University Press, 2017.
    
    \item A.~Rakhlin, K.~Sridharan, and A.B.~Tsybakov. Empirical entropy, minimax regret and minimax risk. \textit{Bernoulli}, 23(2):789--824, 2017.
\end{enumerate}

\end{document}
